% theory.tex -- unified NoT paper. Prose form of papers/unified/theory.md. No lists, no
% em dashes, no colons or semicolons in running prose. Numbers carry % source comments.
\section{Theory}
\label{sec-theory}

A model's answer is pulled toward the interest that is present in its
context. On an open-ended dilemma the interest is the asker's own account
of the conflict, on a false premise it is the premise the user owns, and
on an assigned side it is the side. The pull cannot be removed from inside
the same context, because the copy that would remove it is the copy that
carries it. Three copies of one model vote as one voter
% 16.15.4: majority = neutral seat on 1,642 of 1,642 debates; CV oracle +0.000
and the gain from a correlated jury is bounded by how far its members'
errors are independent
\citep{ladha1992condorcet,boland1989majority,dietrich2013independent,kim2025correlated}.
What can be done is to place an opposed interest into the process so that
agreement across interests, or dissent against interest, becomes
informative. This is the argumentative theory of reasoning
\citep{mercier2011argumentative,mercier2017enigma} joined to the social
account of objectivity \citep{longino1990science,longino2002fate,popper2011open}.
Individuals argue for a side, and truth is a property of the structure of
criticism among them rather than of any one mind. The scaffold and the
collective are the two ways of placing the opposed interest.

Internal opposition is the scaffold. One trace is made to represent the
other parties. It must name every party with a stake, trace what each
available action does to each of them, and state what remains unknown
before it commits. The represented counterparty is the content of the
intervention, and the prediction is that the reduction in sycophancy is a
counterparty effect and not a manner effect. Three consequences were
testable, and they resolve differently. It should be concentrated on
accounts that contain a party other than the asker and absent where none
exists, and it is: the reduction is three times as large on accounts
that name a second party than on those that do not, and a false theorem,
which has no second party, is the one place the scaffold affirms the
falsehood more often rather than less.
% counterparty_moderator.json pooled -37.7 vs -11.3, diff -26.3 [-35.7,-17.0], Z2 RESULTS
% BrokenMath +6.7 [3.3, 10.5]; 16.15.5 S0 +0.067 [+0.033, +0.103]; 16.15.6 abstention reading
It should be a property of the prompt's structure, reproducing across
vendors and on open-weight models, and it is, on seven models from six
vendors, three of them open-weight, with the effect if anything larger
on the open-weight deployments than on any proprietary model.
% Addendum 17.1 RESULTS: 22 to 70 points, 31 to 87 percent
But the prediction that the reduction should be carried specifically by
the sections that name and trace consequences for the other parties,
rather than by the section that licenses hedging, does not hold at the
level of the scaffold's own sections. Removing any one of the three
candidate sections, the stakeholder list, the consequence trace, or the
uncertainty statement, takes back a large share of the reduction on one
model and almost none of it on two others, and pooled across models no
single section clears a quarter of the reduction with a confidence
interval excluding zero.
% Addendum 17.3 RESULTS: pooled NEITHER; haiku shares 0.42-0.55 (mostly excluding 0),
% grok and nano shares near zero with wide intervals including zero
The represented-counterparty account holds at the level of what the
account contains, an item that names a second party loses more validation
than one that does not, and does not hold at the level of which sentence
in the scaffold does the work: no one section is the carrier, on most
models the effect is not attributable to any single part, and where it
is attributable at all it does not cleanly separate the counterparty
account from the hedging account. The theory's item-level prediction is
therefore supported and its section-level prediction is not, a
distinction the paper states rather than collapses.

External opposition is the collective. When the interests are separate
seats the credible-signal rule applies. A statement against the speaker's
assigned interest is informative and a statement with it is not
\citep{milgrom1986relying,lipman1995robust,glazer2001debates}. A verdict
with one losing party gives that party's seat a reason to object, and its
objection carries nothing. A second objection has no side to explain it
and must come from the account. The dissent edge exists only where the
interests are opposed, and it survives removing the exchange between
seats entirely, down to about half its discriminating power rather than
to nothing
\citep{dewatripont1999advocates,krishna2001model}.
% 16.10 fire 0.196 embodied vs 0.005 identical; edges off 0.146; 16.25 no-edge k=4 within-one-loser
% lift 5.06x [+0.114, +0.347] at 78 fired against 9.83x [+0.201, +0.493] on-edge; advocate phi
% -0.842 [-0.876, -0.802] with edges off against -0.901 on-edge (16.12), still strongly opposed
So a copy that never reads another seat's argument, and sees only the
account and the moderator's verdict, still objects almost entirely on the
strength of its assigned interest. The exchange between seats adds the
remaining half, not the phenomenon itself.

The sign of the mechanism is not uniform across models, and the lift is
more robust than the sign. On the model with the highest role-lock the two
advocates' objections are strongly anti-correlated, exactly as the
credible-signal rule predicts: the losing seat objects and the winning
seat does not. On two other models, one with low role-lock and one with
role-lock nearly as high as the first, the two advocates' objections
correlate positively instead, closer to two independent readers both
noticing a hard case than to one side conceding against interest.
% 16.15.1 one-loser phi -0.901 (grok); 16.23 RESULTS phi +0.246 (haiku, lock 0.419);
% 16.24 RESULTS phi +0.148 [+0.065, +0.227] (Llama, lock 0.862, high but not opposed)
The within-one-loser accuracy lift is positive with a confidence interval
excluding zero in every one of these conditions regardless of the sign of
the correlation, so role-lock alone does not decide whether the
against-interest mechanism specifically is operating, only whether the
lift is present at all. The theory's claim that objectivity comes from
placing an opposed interest into the process is supported by the lift's
robustness; the specific reading rule, that a second objection is
informative because it goes against the objecting seat's own assigned
side, is confirmed on the model that locks to its role and is not
established on the others, where the lift more plausibly reflects
correlated difficulty than opposed interest.

The collective is a sensor and not a decider. No arrangement beats its best
seat by vote, the moderator's competence sets the verdict
\citep{shin1998adversarial},
% 16.17 sonnet +0.038 [+0.014, +0.064], haiku -0.115 [-0.169, -0.060]
and the gain is realised by routing the flagged cases to a stronger judge
with no side. This part of the claim is the most robust of all: routing
beats the collective alone on every condition tested, on-edge and
off-edge, on the model that locks to its role and on the two that do not,
independently reproduced by a second implementation in every case.
% verify_pillar3_headline.py: grok +0.070 [+0.048, +0.094], no-edge k=4 +0.048 [+0.032, +0.065],
% haiku-249 +0.037 [+0.015, +0.059], Llama +0.041 [+0.028, +0.056]; none has failed
Beating the strongest single model, rather than the collective, is
weaker and model-specific: it holds on the role-locking model and does
not on the others.
% 16.22 R1 grok +0.024 [+0.005, +0.045]; 16.24 Llama R1 CI includes 0
That is the precise sense in which the collective depends on the same
control the scaffold exploits. Embodiment has to hold for the
against-interest reading of opposition to be readable, even though a
weaker, sign-agnostic dissent signal survives without it.

The game-theoretic reading is used only where its primitives exist, an
assigned interest, a receiver whose verdict is separate from the message,
and a decision read independently of the speaker. Those exist in the
collective and not in a single trace, so for the scaffold alone the paper
claims a change in the output distribution whose content is the
represented counterparty and makes no signalling claim.

\todo{Pillar two, theory of mind (represented minds). Naming what another
party knows and wants is an expression-side act, and whether the scaffold
closes the gap between what a model represents about other minds and what
it expresses is a live question this paper does not answer. Out of scope
for the present work (PI decision, 2026-09-21); left here as a stub for a
future developer. Stage-1 data exists (Addendum 18, four models, 320
ToMBench items, a guard-clean negative result with an identified
mechanism) but is not part of this paper's claims and is not pooled.}

This paper validates two of the arrangement's readouts. The asker's pull
on a single trace is pillar one, and the readable dissent that separate
seats produce is pillar three. Two results tie the pillars to each other
rather than merely placing them side by side. Narration does not create
the collective's structure, so internal and external opposition are
different objects with different evidence. And without role lock the
external structure collapses to correlated readers, so external opposition
rests on the control that assignment gives over what a copy will say.
